\customchapter{ABSTRACT}

\begin{center}
\large \vspace{-1.5cm} \textbf{HyperNCD--Sonata hybrid model for novel class
discovery in the semantic segmentation of architectural heritage point clouds}
\end{center}

The digital recording of architectural heritage produces point clouds that show
where matter is, but not which element each point stands for. Three-dimensional
semantic segmentation methods work over a closed list of categories and assign
every unforeseen element to the closest known class. The objective is to
measurably improve the recognition of elements absent from the training labels.

This is applied, quantitative research with a comparative experimental design. We
first diagnose the best available method for novel class discovery through its
ablation study, and then replace its weakest component with a
self-supervised encoder. The evaluation uses the survey of the Cathedral of Lima,
two gigabytes in five sections, against a three-run baseline.

The base method reaches between 11.8 and 47.5 points of intersection over union
on the unseen categories, against 49.8 for the supervised equivalent. Its
ablation credits the geometric descriptor with only 4.2 points over a base of
23.2, less than half of the 9.6 of the hypergraph. Removing that dependence on
low-level signals raises a frozen descriptor from 5.6 to 72.5 points. We also
define a closed-gap measure, since correcting the bias towards majority classes
adds up to 15.34 points.

The imbalance is documented by the authors themselves and remains uncorrected,
but whether it persists without labels for the target categories is still to be
established. Every point gained is correction work the specialist no longer does.

\vspace{1em}

\noindent\begin{minipage}{\linewidth}
\noindent\textbf{Keywords:}\par
\noindent Point clouds; Three-dimensional semantic segmentation; Novel class
discovery; Self-supervised learning; Hypergraphs; Architectural heritage
\end{minipage}
